\documentclass[11pt,a4paper]{article}

% --- Essential Packages ---
\usepackage[utf8]{inputenc}
\usepackage[margin=1in]{geometry}
\usepackage{amsmath,amssymb,amsfonts}
\usepackage{booktabs}
\usepackage{tabularx}
\usepackage{array}
\usepackage{xcolor}
\usepackage{hyperref}
\usepackage{enumitem}
\usepackage{caption}
\usepackage{microtype}
\usepackage{fancyhdr}

% Allow line breaks at underscores in texttt
\usepackage[hyphens]{url}
\newcommand{\code}[1]{\texttt{\seqsplit{#1}}}
% If seqsplit is not present in standard, use breakable underscore:
\renewcommand{\_}{\textunderscore\allowbreak}

% --- Colors & Styling ---
\definecolor{primary}{RGB}{24, 43, 73}
\definecolor{secondary}{RGB}{70, 130, 180}
\definecolor{darkgreen}{RGB}{34, 139, 34}

\hypersetup{
    colorlinks=true,
    linkcolor=primary,
    citecolor=secondary,
    urlcolor=secondary
}

\pagestyle{fancy}
\fancyhf{}
\rhead{\small \textcolor{gray}{Collusion Detection Multi-Detector Benchmark Report}}
\lhead{\small \textcolor{gray}{Procurement Auction Analytics}}
\rfoot{\small Page \thepage}

\title{\textbf{\Large Empirical Benchmark \& Comparative Analysis of Collusion Detectors:}\\
\large Testing 29 Systematic Metric Combinations Across Real-World Procurement Cartels}
\author{\textbf{Research \& Evaluation Report} \\ 
Department of Computer Science / Economics \\
Detection of Collusion in Procurement Auctions Project}
\date{\today}

\begin{document}

\maketitle
\thispagestyle{empty}

\begin{abstract}
Screening for cartels and bid-rigging in public procurement is complicated by institutional heterogeneity across auction formats and the disparate tactics employed by colluding firms (e.g., bid rotation, cover bidding, and market allocation). This report provides a comprehensive empirical evaluation of \textbf{29 systematically constructed collusion detectors} across five structural tiers: Single-Metric baselines, Pairwise metrics, Triplet metrics, Quad metrics (leave-one-group-out ablations), and an All-Metrics ensemble. We test these architectures on authentic procurement auction cartels from \textbf{Japan (Okinawa)}, \textbf{Italy}, and \textbf{Brazil} (1,459 tenders) using a rigorous \textbf{Leave-One-Dataset-Out (LODO)} cross-validation protocol. Our primary finding is that the All-Metrics model is sub-optimal: the global winner is \texttt{real\_quad\_minus\_corr\_gbm} (combining price dispersion, distribution moments, winning gaps, and bidder-pair cooperation while ablating raw residual correlation), achieving a Macro Real AUC of \textbf{0.745} and Precision@10\% of \textbf{0.648}, compared to \textbf{0.726} AUC for the All-Metrics detector and $\le \mathbf{0.656}$ for single-metric screens.
\end{abstract}

\vspace{0.5em}
\hrule
\vspace{1.5em}

\section{Introduction \& Problem Formulation}

Antitrust authorities face severe challenges when screening procurement data for bid-rigging. While individual statistical screens (such as the coefficient of variation or normalized winning gaps) have been proposed in empirical industrial organization literature, cartels frequently exploit structural loopholes or adapt their bidding patterns to mimic competitive distributions. 

A natural hypothesis is that aggregating all available statistical screens into a machine learning classifier will maximize detection accuracy. However, combining heterogeneous metrics across diverse regulatory jurisdictions risks introducing multicollinearity, spurious correlations, and out-of-domain overfitting.

To rigorously evaluate this trade-off, we implemented a \textbf{Multi-Detector Framework} that systematically initializes detectors across five architectural tiers:
\begin{enumerate}[noitemsep]
    \item \textbf{Single-Metric Detectors (Tier 1)}: Isolated statistical screens testing individual behavioral signals.
    \item \textbf{Pairwise Metric Detectors (Tier 2)}: Two-group combinations pairing complementary dimensions (e.g., price dispersion + firm identity tracking).
    \item \textbf{Triplet Metric Detectors (Tier 3)}: Three-group structural combinations.
    \item \textbf{Quad Metric Detectors (Tier 4)}: Four-group models implementing systematic leave-one-group-out feature ablations.
    \item \textbf{All-Metrics Ensemble (Tier 5)}: A comprehensive detector utilizing every defined metric group simultaneously.
\end{enumerate}

---

\section{Metric Taxonomy \& Mathematical Formulations}

We evaluate collusion signals across two complementary environments:
\begin{enumerate}[noitemsep]
    \item \textbf{Longitudinal Market Level (Simulated Framework)}: 53 core features capturing market-wide dynamic coordination across sequential tenders.
    \item \textbf{Individual Tender Level (Empirical Real Data)}: 15 features across 5 logical groups capturing cross-national tender anomalies in authentic cartels.
\end{enumerate}

\subsection{Empirical Metric Groups on Real Data (\texttt{BIDRIG\_REAL\_DATA})}

The 15 real-data features are partitioned into 5 logical groups based on their behavioral and economic mechanisms (Table~\ref{tab:real_groups}).

\begin{table}[h!]
\centering
\small
\caption{Empirical Metric Groups in Real Procurement Data}
\label{tab:real_groups}
\begin{tabularx}{\textwidth}{lp{4.2cm}X}
\toprule
\textbf{Group} & \textbf{Features Included} & \textbf{Mathematical Definition \& Economic Mechanism} \\
\midrule
\textbf{disp} & \texttt{scr\_cv}, \texttt{scr\_spread}, \texttt{scr\_log\_nbids} & 
\textbf{Dispersion \& Turnout}: $\text{CV} = \sigma_b / \mu_b$, spread $(b_{\max} - b_{(1)})/b_{(1)}$, and log bidder count $\log(N)$. Cartels suppress variance to guarantee stable profit division. \\
\addlinespace
\textbf{shape} & \texttt{scr\_skew}, \texttt{scr\_kurt} & 
\textbf{Higher Distribution Moments}: Standardized skewness and kurtosis ($\frac{m_4}{m_2^2} - 3$). Extreme values indicate focal pricing or unnatural clustering. \\
\addlinespace
\textbf{gap} & \texttt{scr\_rd}, \texttt{scr\_rdalt}, \texttt{scr\_gap} & 
\textbf{Winning Margins}: Normalized winning gap $\frac{b_{(2)} - b_{(1)}}{b_{(1)}}$ and Relative Distance ($\text{RD} = \frac{b_{(2)} - b_{(1)}}{\text{SD}(\text{losing bids})}$). Large gaps expose deliberate cover bidding. \\
\addlinespace
\textbf{pair\_corr} & \texttt{pair\_corr\_max}, \texttt{pair\_corr\_top3}, \texttt{pair\_coverage} & 
\textbf{Historical Residual Correlation}: Max and top-3 Fisher-$z$ transformed residual correlations computed across all other tenders in that jurisdiction. \\
\addlinespace
\textbf{pair\_coop} & \texttt{pair\_close\_max}, \texttt{pair\_close\_top3}, \texttt{pair\_cobid\_max}, \texttt{pair\_cobid\_top3} & 
\textbf{Co-bidding Cooperation}: Share of historical tenders where pairs bid within 3\% ($\|b_i - b_j\|/b_i \le 0.03$) and their co-bidding correlation coefficient $\phi$. \\
\bottomrule
\end{tabularx}
\end{table}

\subsection{Simulated Benchmark Metric Groups}
For multi-tender simulated markets, the feature set expands to 53 features across 5 groups:
\begin{itemize}[noitemsep]
    \item \textbf{\texttt{lit}} (21 features): Per-tender screens (\texttt{cv}, \texttt{skew}, \texttt{kurt}, \texttt{rd}, \texttt{rdalt}, \texttt{gap}, \texttt{spread}) aggregated over the market via mean, median, and standard deviation.
    \item \textbf{\texttt{dep}} (22 features): Identity-dependence statistics calibrated against both within-tender permutation nulls and time-shuffled firm residual nulls.
    \item \textbf{\texttt{part}} (3 features): Exclusivity indices, including greedy search for cliques of firms that avoid co-bidding (\texttt{part\_group}).
    \item \textbf{\texttt{change}} (4 features): Brownian-bridge partial sum scans detecting abrupt structural breaks in variance and winning gaps.
    \item \textbf{\texttt{seq}} (3 features): Dynamic Page-Hinkley CUSUM and 2-state Hidden Markov Model forward filter posteriors $P(\text{State} = \text{Collusive} \mid u_{1:t})$.
\end{itemize}

---

\section{Empirical Evaluation Methodology}

We evaluated all 29 detector combinations against authentic procurement cartels from three distinct jurisdictions:
\begin{itemize}[noitemsep]
    \item \textbf{Japan (Okinawa)}: 1,080 tenders (123 collusive, 1,665 persistent firms, First-Price with reserve price floor).
    \item \textbf{Italy}: 278 tenders (143 collusive, 1,029 firms, Average-Bid mechanism where the bid closest to the trimmed mean wins).
    \item \textbf{Brazil}: 101 tenders (33 collusive, 247 firms, First-Price procurement tenders with sparse bidder overlap).
\end{itemize}

\subsection{Validation Protocol}
\begin{enumerate}[noitemsep]
    \item \textbf{Leave-One-Dataset-Out (LODO)}: Models train on two countries and evaluate strictly out-of-domain on the third held-out country.
    \item \textbf{Percentile Scaling}: Features are transformed to within-dataset percentiles (label-free) to ensure cross-currency comparability.
    \item \textbf{False Alarm Calibration}: The classification threshold is fixed at a 5\% False Alarm Rate (FAR) on out-of-fold competitive tenders.
    \item \textbf{Autocorrelation Robustness}: Moving-block bootstrapping (block length 20 tenders, 300 resamples) yields robust 95\% confidence intervals.
\end{enumerate}

---

\section{Empirical Results \& Comparative Analysis}

To provide complete clarity without overflowing the printed page width, the results are divided into two focused tables:
\begin{itemize}[noitemsep]
    \item \textbf{Table~\ref{tab:macro_perf}}: Macro Out-of-Domain Performance (Ranking, Metric Composition, Number of Features, Macro AUC, and Macro Precision@10\%).
    \item \textbf{Table~\ref{tab:jurisdiction_breakdown}}: Jurisdiction-by-Jurisdiction Generalization (Japan, Italy, Brazil AUCs, and False Alarm Rates).
\end{itemize}

\begin{table}[h!]
\centering
\small
\caption{Overall Out-of-Domain Benchmark Performance: Single Metrics vs. Top 5 Combinations vs. All-Metrics Ensemble}
\label{tab:macro_perf}
\begin{tabularx}{\textwidth}{llXrrr}
\toprule
\textbf{Tier / Category} & \textbf{Detector Identifier} & \textbf{Metric Groups Used} & \textbf{\# Feat} & \textbf{Macro AUC} & \textbf{Prec@10\%} \\
\midrule
\multicolumn{6}{l}{\textit{\textbf{Top 5 Performing Combinations}}} \\
\textbf{Rank 1 (Winner)} & \texttt{real\_quad\_minus\_corr\_gbm} & \texttt{disp + shape + gap + pair\_coop} & 12 & \textbf{0.745} & \textbf{0.648} \\
\textbf{Rank 2} & \texttt{real\_triplet\_disp\_shape\_coop} & \texttt{disp + shape + pair\_coop} & 9 & \textbf{0.735} & 0.509 \\
\textbf{Rank 3} & \texttt{real\_quad\_minus\_gap\_gbm} & \texttt{disp + shape + pair\_corr + pair\_coop} & 12 & \textbf{0.731} & 0.545 \\
\textbf{Rank 4} & \texttt{real\_triplet\_disp\_corr\_coop} & \texttt{disp + pair\_corr + pair\_coop} & 10 & \textbf{0.725} & 0.499 \\
\textbf{Rank 5} & \texttt{real\_quad\_minus\_shape\_gbm} & \texttt{disp + gap + pair\_corr + pair\_coop} & 13 & \textbf{0.723} & 0.548 \\
\midrule
\multicolumn{6}{l}{\textit{\textbf{All-Metrics Ensemble}}} \\
\textbf{Benchmark} & \texttt{real\_all\_metrics\_gbm} & \texttt{All 5 Groups Combined} & 15 & \textbf{0.726} & 0.538 \\
\midrule
\multicolumn{6}{l}{\textit{\textbf{Single-Metric Baselines}}} \\
Baseline 1 & \texttt{real\_disp\_gbm} & \texttt{disp} & 3 & \textbf{0.656} & 0.374 \\
Baseline 2 & \texttt{real\_coop\_gbm} & \texttt{pair\_coop} & 4 & \textbf{0.626} & 0.444 \\
Baseline 3 & \texttt{real\_gap\_gbm} & \texttt{gap} & 3 & \textbf{0.608} & 0.364 \\
Baseline 4 & \texttt{real\_shape\_gbm} & \texttt{shape} & 2 & \textbf{0.545} & 0.372 \\
Baseline 5 & \texttt{real\_corr\_gbm} & \texttt{pair\_corr} & 3 & \textbf{0.540} & 0.332 \\
\bottomrule
\end{tabularx}
\end{table}

\begin{table}[h!]
\centering
\small
\caption{Cross-Jurisdiction Performance Breakdown Across Regulatory Formats}
\label{tab:jurisdiction_breakdown}
\begin{tabularx}{\textwidth}{lXXXXr}
\toprule
\textbf{Detector Identifier} & \textbf{Japan AUC} \newline \footnotesize{(Floor Reserve)} & \textbf{Italy AUC} \newline \footnotesize{(Average-Bid)} & \textbf{Brazil AUC} \newline \footnotesize{(First-Price)} & \textbf{Macro AUC} & \textbf{False Alarm Rate} \\
\midrule
\multicolumn{6}{l}{\textit{\textbf{Top 5 Performing Combinations}}} \\
\texttt{real\_quad\_minus\_corr\_gbm} & 0.690 & 0.765 & \textbf{0.780} & \textbf{0.745} & 0.102 \\
\texttt{real\_triplet\_disp\_shape\_coop} & 0.664 & 0.782 & 0.759 & \textbf{0.735} & 0.101 \\
\texttt{real\_quad\_minus\_gap\_gbm} & 0.670 & \textbf{0.783} & 0.740 & \textbf{0.731} & 0.138 \\
\texttt{real\_triplet\_disp\_corr\_coop} & 0.679 & 0.749 & 0.746 & \textbf{0.725} & 0.155 \\
\texttt{real\_quad\_minus\_shape\_gbm} & \textbf{0.692} & 0.743 & 0.735 & \textbf{0.723} & 0.134 \\
\midrule
\multicolumn{6}{l}{\textit{\textbf{All-Metrics Ensemble}}} \\
\texttt{real\_all\_metrics\_gbm} & 0.687 & 0.757 & 0.734 & \textbf{0.726} & 0.129 \\
\midrule
\multicolumn{6}{l}{\textit{\textbf{Single-Metric Baselines}}} \\
\texttt{real\_disp\_gbm} & 0.602 & 0.653 & 0.715 & \textbf{0.656} & 0.198 \\
\texttt{real\_coop\_gbm} & 0.649 & 0.640 & 0.590 & \textbf{0.626} & 0.110 \\
\texttt{real\_gap\_gbm} & 0.617 & 0.509 & 0.698 & \textbf{0.608} & 0.079 \\
\texttt{real\_shape\_gbm} & 0.456 & 0.568 & 0.611 & \textbf{0.545} & 0.102 \\
\texttt{real\_corr\_gbm} & 0.547 & 0.544 & 0.530 & \textbf{0.540} & 0.063 \\
\bottomrule
\end{tabularx}
\end{table}

---

\section{Discussion \& Key Findings}

\subsection{1. The All-Metrics Dilemma: The Cost of Overfitting}
The empirical results definitively demonstrate that \textbf{including all metrics is not optimal for cross-domain cartel detection}. The All-Metrics detector (\texttt{real\_all\_metrics\_gbm}, Macro AUC 0.726) is outperformed by three smaller combinations:
\begin{itemize}[noitemsep]
    \item \texttt{real\_quad\_minus\_corr\_gbm} (12 features) achieves Macro AUC \textbf{0.745} (+0.019 gain) and Precision@10\% of \textbf{0.648} (+0.110 gain).
    \item \texttt{real\_triplet\_disp\_shape\_coop\_gbm} (9 features) achieves Macro AUC \textbf{0.735}.
\end{itemize}
\textbf{Theoretical Rationale}: Raw pairwise residual correlation (\texttt{pair\_corr}) is highly vulnerable to dataset-specific idiosyncrasies. When transferring across jurisdictions where firm populations do not overlap, \texttt{pair\_corr} introduces noise that degrades tree splits. In contrast, behavioral co-bidding cooperation (\texttt{pair\_coop}) generalizes effectively across institutional rules.

\subsection{2. Failure Modes of Single-Metric Screens}
Deploying statistical screens in isolation produces fragile detectors:
\begin{itemize}[noitemsep]
    \item While \texttt{real\_disp\_gbm} performs decently on Brazil (0.715 AUC), its overall macro AUC is only \textbf{0.656}.
    \item Distribution moments (\texttt{real\_shape\_gbm}, AUC 0.545) and pair correlation (\texttt{real\_corr\_gbm}, AUC 0.540) perform barely above random chance when unassisted.
    \item \textbf{Synergistic Effect}: Pairing price dispersion (\texttt{disp}) with bidder cooperation (\texttt{pair\_coop}) immediately lifts Macro AUC from 0.656 to \textbf{0.707}, and adding distribution shape reaches \textbf{0.735}.
\end{itemize}

\subsection{3. Institutional Heterogeneity Across Jurisdictions}
The performance variations highlight the impact of auction design:
\begin{itemize}[noitemsep]
    \item \textbf{Italy (Average-Bid Mechanism)}: Because the winner is the firm closest to the average bid, lowest-bid screens break down completely (\texttt{real\_gap\_gbm} yields an AUC of only 0.509). However, \texttt{pair\_coop} and \texttt{pair\_corr} excel, driving \texttt{real\_quad\_minus\_gap\_gbm} to an outstanding \textbf{0.783} AUC.
    \item \textbf{Brazil (First-Price, Sparse Bidders)}: With minimal historical pair overlap, distribution screens dominate (\texttt{disp+shape} achieves 0.769 AUC).
    \item \textbf{Japan (High Firm Persistence)}: High bidder recurrence allows gap and co-bidding metrics to excel (\texttt{real\_pair\_gap\_coop\_gbm} achieves 0.717 AUC).
\end{itemize}

---

\section{Conclusion \& Practical Recommendations}

For antitrust regulators and procurement oversight agencies:
\begin{enumerate}[noitemsep]
    \item \textbf{Do Not Rely on Single Screens}: Isolated metrics suffer high false alarm rates and complete blindness under alternative auction rules.
    \item \textbf{Avoid Unfiltered Feature Bloat}: Pruning noisy, dataset-specific correlation features in favor of behavioral co-bidding cooperation yields superior cross-market generalization.
    \item \textbf{Adopt the Quad Architecture}: The \texttt{real\_quad\_minus\_corr\_gbm} configuration provides the most robust, transferable defense against procurement collusion worldwide.
\end{enumerate}

\end{document}
